\subsection{形式光滑代数的例子}

设 k 为环。

\begin{enumerate}
\def\labelenumi{\arabic{enumi})}
\tightlist
\item
  设 P 为投射 \(k\)-模。对称 \(k\)-代数 \(\mathbf{S}_k(\mathrm{P})\) 对离散拓扑形式光滑，更不必说对其分次所定义的拓扑形式光滑。事实上，对任意 \(k\)-代数 C 与 C 的每个理想 N，从 \(\mathbf{S}_k(\mathrm{P})\) 到 C（相应地 C/N）中的代数同态与从 P 到 C（相应地 C/N）中的 \(k\)-线性映射一一对应，且典范映射 \(\mathrm{Hom}_k(\mathrm{P},\mathrm{C}) \to \mathrm{Hom}_k(\mathrm{P},\mathrm{C}/\mathrm{N})\) 是满射的。
\end{enumerate}

因此 (命题 3, c))，\(k\)-代数 \(\widehat{\mathbf{S}}_k(\mathrm{P}) = \prod_{n\geqslant 0}\mathbf{S}_k^n (\mathrm{P})\) 是形式光滑的（对诸 \(\mathbf{S}_k^n (\mathrm{P}))\) 上离散拓扑的乘积拓扑）；事实上，它是 \(k\)-代数 \(\mathbf{S}_k(\mathrm{P})\) 对由分次定义的拓扑的完备化。

\begin{enumerate}
\def\labelenumi{\arabic{enumi})}
\setcounter{enumi}{1}
\item
  对任意不定元族 \(\mathbf{T} = (\mathrm{T}_i)_{i\in I}\)，多项式 \(k\)-代数 \(k[\mathbf{T}]\) 与赋予其典范拓扑的形式幂级数 \(k\)-代数 \(k[[\mathbf{T}]]\) 都是形式光滑的；这由例 1 得出。若 \(k\) 是域，则纯超越扩张 \(k(\mathbf{T})\) 形式光滑 (第 2 节, 命题 4 a))。
\item
  设 \(f \in k[\mathrm{T}]\) 为一个不定元的多项式。说 \(k\)-代数 \(k[\mathrm{T}]/(f)\) 形式光滑，就是说满足下述性质：对每个 \(k\)-代数 C 与 C 的每个平方零理想 N，\(f\) 在 C/N 中的每个根都可提升为 \(f\) 在 C 中的根。当 \(f\) 与其导数 \(f'\) 生成单位理想时即属此情形。事实上，设 \(\alpha\) 为 \(f\) 在 C/N 中的根，\(a\) 为 C 中提升 \(\alpha\) 的元素。于是 \(f(a)\) 属于 N，从而 \(f'(a)\) 在 C 中可逆；元素 \(b = a - f'(a)^{-1}f(a)\) 提升 \(\alpha\)。由于 \(f'(a)^{-1}f(a)\) 的平方为零，有
\end{enumerate}

\[
f (b) = f (a) - f ^ {\prime} (a) f ^ {\prime} (a) ^ {- 1} f (a) = 0.
\]

定理 1 (I. S. Cohen). 设 \(k\) 为域，\(K\) 为 \(k\) 的可分扩张。则 \(K\) 是形式光滑 \(k\)-代数。

设 C 为 \(k\)-代数，N 为 C 中的平方零理想，\(\pi: \mathrm{C} \to \mathrm{C/N}\) 为典范同态，\(\varphi: \mathrm{K} \to \mathrm{C/N}\) 为 \(k\)-代数同态。尚需构造 \(\varphi\) 的一个提升。我们区分两种情形。

\begin{enumerate}
\def\labelenumi{\Alph{enumi})}
\item
  先设 k 的特征为 0。考察偶 \((\mathrm{K}^{\prime}, \tilde{\varphi}^{\prime})\)，其中 \(K^{\prime}\) 是 K 的中间扩张，\(\tilde{\varphi}^{\prime}: K^{\prime} \to C\) 是 \(\varphi\) 在 \(K^{\prime}\) 上的限制的一个提升。这些偶按延拓关系排序后的集合是归纳的；由 Zorn 引理 (E, III, 第 20 页, 定理 2)，存在极大偶 \((\mathrm{K}^{\prime}, \tilde{\varphi}^{\prime})\)。我们来证明 \(K^{\prime}\) 等于 K。设 \(x \in K - K^{\prime}\)。若 x 在 \(K^{\prime}\) 上超越，则 \(K^{\prime}\)-代数 \(\mathrm{K}^{\prime}(x)\) 形式光滑 (例 2)。若 x 在 \(K^{\prime}\) 上代数，则其极小多项式 \(f \in K^{\prime}[T]\) 与其导数互素 (A, V, 第 37 页, 命题 4)，而 \(\mathrm{K}^{\prime}(x)\) 与 \(K^{\prime}\)-代数 \(\mathrm{K}^{\prime}[\mathrm{T}]/(f)\) 恒同，后者因此是形式光滑 \(K^{\prime}\)-代数 (例 3)。两种情形下，\(\mathrm{K}^{\prime}(x)\) 都在 \(K^{\prime}\) 上形式光滑，于是存在 \(\tilde{\varphi}^{\prime}\) 到 \(\mathrm{K}^{\prime}(x)\) 的一个延拓，它提升 \(\varphi\) 在 \(\mathrm{K}^{\prime}(x)\) 上的限制，这与 \((\mathrm{K}^{\prime}, \tilde{\varphi}^{\prime})\) 的极大性矛盾。
\item
  设 \(k\) 的特征为 \(p \neq 0\)。考察环同态 F: C → C，它满足 F(x) = x \(^{p}\)；对 x ∈ N 有 F(x) = 0，于是存在唯一的环同态 λ: C/N → C 使 λ ∘ π = F。有 π(λ(π(x))) = π(x \(^{p}\)) = π(x) \(^{p}\)；由于 π 满射，故对 C/N 的每个元素 z 有 π(λ(z)) = z \(^{p}\)。另一方面，设 f: K → K \(^{p}\) 为同构 y ↦ y \(^{p}\)，f \(^{-1}\): K \(^{p}\) → K 为其逆同构。设 g: K \(^{p}\) → C 为下列环同态的复合
\end{enumerate}

\[
\mathrm{K} ^ {p} \xrightarrow {f ^ {- 1}} \mathrm{K} \xrightarrow {\varphi} \mathrm{C} / \mathrm{N} \xrightarrow {\lambda} \mathrm{C}.
\]

对每个 \(x \in K\)，有 \(g(x^{p}) = \lambda(\varphi(x))\)。由于 \(\lambda(\alpha z) = \alpha^{p}\lambda(z)\)（对 \(\alpha \in k\) 与 \(z \in C/N\)），映射 \(g\) 是 \(k^{p}\)-线性的。由于 \(K\) 在 \(k\) 上的扩张可分，\(k(K^{p})\) 与 \(k \otimes_{k^{p}} K^{p}\) 恒同 (A, V, 第 119 页, 注)；于是存在唯一的 \(k\)-代数同态 \(h: k(K^{p}) \to C\)，它与 \(g\) 在 \(K^{p}\) 上重合。

设 \((a_{i})_{i\in I}\) 为 K 的一个 \(p\)-基（在 \(k(\mathrm{K}^p)\) 上）(A, V, 第 98 页, 定理 2)；对每个 \(i\in I\)，选取 C 的一个元素 \(b_{i}\)，满足 \(\pi(b_{i})=\varphi(a_{i})\)。有 \(h(a_{i}^{p})=g(a_{i}^{p})=\lambda(\varphi(a_{i}))=\lambda(\pi(b_{i}))=b_{i}^{p}\)（对每个 \(i\in I\)）。据 A, V, 第 94 页的注，存在 \(k\)-代数同态 \(\tilde{\varphi}:K\to C\)，它扩张 \(h\) 且满足 \(\tilde{\varphi}(a_{i})=b_{i}\)（对每个 \(i\)）。有 \(\pi(\tilde{\varphi}(a_{i}))=\pi(b_{i})=\varphi(a_{i})\)（对每个 \(i\)），且 \(\pi(\tilde{\varphi}(x^{p}))=\pi(h(x^{p}))=\pi(g(x^{p}))=\pi(\lambda(\varphi(x)))=\varphi(x^{p})\)（对每个 \(x\in K\)）。于是有 \(\pi\circ\tilde{\varphi}=\varphi\)，证明完毕。

推论.--- 设 k 为域，K 为 k 的可分扩张，A 为线性拓扑 K-代数。若 A 在 K 上形式光滑，则它在 k 上形式光滑。

这由定理及第 2 节的命题 3 a) 得出。

注.- 1) 设 \(k\) 为域。每个平展的 \(k\) -代数 (A, V, 第 28 页, 定义 1) 都形式光滑 (同上引文, 第 34 页, 定理 4, d) 与第 2 节, 命题 3, b))。

\begin{enumerate}
\def\labelenumi{\arabic{enumi})}
\setcounter{enumi}{1}
\tightlist
\item
  我们将在下面 (第 5 节定理 2 的推论 2) 看到：形式光滑的域扩张是绝对正则的，从而是可分的 (§ 6, 第 4 节, 例 2)。
\end{enumerate}

